\def\tituloconcepto{VQE}
% Si el capitulo se compila solo (sin main.tex), esta cabecera abre el documento.
% Antes de % Si el capitulo se compila solo (sin main.tex), esta cabecera abre el documento.
% Antes de % Si el capitulo se compila solo (sin main.tex), esta cabecera abre el documento.
% Antes de \input{cabecera} cada capitulo define \tituloconcepto.
\ifdefined\enmain\else
  \documentclass[11pt,a4paper]{article}
  \input{preambulo}
  \begin{document}
  \title{\tituloconcepto}
  \author{\autor}
  \date{\today}
  \maketitle
\fi
 cada capitulo define \tituloconcepto.
\ifdefined\enmain\else
  \documentclass[11pt,a4paper]{article}
  % Preambulo compartido por todos los documentos. Compilar con XeLaTeX/LuaLaTeX
% (o Tectonic). Si se usa Overleaf, elegir el compilador XeLaTeX.
\usepackage[spanish,es-noshorthands]{babel}
\usepackage{amsmath,amssymb}
\usepackage{booktabs,array}
\usepackage[a4paper,margin=2.5cm]{geometry}
\usepackage{xcolor}
\usepackage{listings}
\usepackage[most]{tcolorbox}
\usepackage{quantikz}
\usepackage[hidelinks]{hyperref}

% ---- Autor (editar a gusto) -------------------------------------------------
\newcommand{\autor}{Genesis Silesky}

% ---- Notacion de Dirac y atajos ---------------------------------------------
% quantikz ya define \ket y \bra; DeclareRobustCommand los redefine sin error.
\DeclareRobustCommand{\ket}[1]{\left|#1\right\rangle}
\DeclareRobustCommand{\bra}[1]{\left\langle#1\right|}
\DeclareRobustCommand{\braket}[2]{\left\langle#1\middle|#2\right\rangle}
\newcommand{\C}{\mathbb{C}}
\newcommand{\R}{\mathbb{R}}
\DeclareMathOperator{\Tr}{Tr}

% ---- Cajas de apoyo --------------------------------------------------------
% "concepto": definiciones previas que hacen falta para entender el capitulo.
\newtcolorbox{concepto}[1]{breakable,colback=blue!4,colframe=blue!55!black,
  fonttitle=\bfseries,title={Concepto previo: #1}}
% "nota": aclaraciones y advertencias.
\newtcolorbox{nota}[1]{breakable,colback=orange!6,colframe=orange!70!black,
  fonttitle=\bfseries,title={#1}}

% ---- Estilos de listados ----------------------------------------------------
\lstdefinestyle{py}{language=Python,basicstyle=\ttfamily\footnotesize,
  breaklines=true,postbreak=\mbox{\textcolor{gray}{$\hookrightarrow$}\space},
  numbers=left,numberstyle=\tiny\color{gray},frame=single,
  commentstyle=\color{green!40!black}\itshape,
  keywordstyle=\color{blue!70!black}\bfseries,
  stringstyle=\color{red!60!black},showstringspaces=false,tabsize=4,
  columns=fullflexible,keepspaces=true}
\lstdefinestyle{salida}{basicstyle=\ttfamily\footnotesize,breaklines=true,
  frame=leftline,backgroundcolor=\color{gray!8},columns=fullflexible,
  keepspaces=true}

  \begin{document}
  \title{\tituloconcepto}
  \author{\autor}
  \date{\today}
  \maketitle
\fi
 cada capitulo define \tituloconcepto.
\ifdefined\enmain\else
  \documentclass[11pt,a4paper]{article}
  % Preambulo compartido por todos los documentos. Compilar con XeLaTeX/LuaLaTeX
% (o Tectonic). Si se usa Overleaf, elegir el compilador XeLaTeX.
\usepackage[spanish,es-noshorthands]{babel}
\usepackage{amsmath,amssymb}
\usepackage{booktabs,array}
\usepackage[a4paper,margin=2.5cm]{geometry}
\usepackage{xcolor}
\usepackage{listings}
\usepackage[most]{tcolorbox}
\usepackage{quantikz}
\usepackage[hidelinks]{hyperref}

% ---- Autor (editar a gusto) -------------------------------------------------
\newcommand{\autor}{Genesis Silesky}

% ---- Notacion de Dirac y atajos ---------------------------------------------
% quantikz ya define \ket y \bra; DeclareRobustCommand los redefine sin error.
\DeclareRobustCommand{\ket}[1]{\left|#1\right\rangle}
\DeclareRobustCommand{\bra}[1]{\left\langle#1\right|}
\DeclareRobustCommand{\braket}[2]{\left\langle#1\middle|#2\right\rangle}
\newcommand{\C}{\mathbb{C}}
\newcommand{\R}{\mathbb{R}}
\DeclareMathOperator{\Tr}{Tr}

% ---- Cajas de apoyo --------------------------------------------------------
% "concepto": definiciones previas que hacen falta para entender el capitulo.
\newtcolorbox{concepto}[1]{breakable,colback=blue!4,colframe=blue!55!black,
  fonttitle=\bfseries,title={Concepto previo: #1}}
% "nota": aclaraciones y advertencias.
\newtcolorbox{nota}[1]{breakable,colback=orange!6,colframe=orange!70!black,
  fonttitle=\bfseries,title={#1}}

% ---- Estilos de listados ----------------------------------------------------
\lstdefinestyle{py}{language=Python,basicstyle=\ttfamily\footnotesize,
  breaklines=true,postbreak=\mbox{\textcolor{gray}{$\hookrightarrow$}\space},
  numbers=left,numberstyle=\tiny\color{gray},frame=single,
  commentstyle=\color{green!40!black}\itshape,
  keywordstyle=\color{blue!70!black}\bfseries,
  stringstyle=\color{red!60!black},showstringspaces=false,tabsize=4,
  columns=fullflexible,keepspaces=true}
\lstdefinestyle{salida}{basicstyle=\ttfamily\footnotesize,breaklines=true,
  frame=leftline,backgroundcolor=\color{gray!8},columns=fullflexible,
  keepspaces=true}

  \begin{document}
  \title{\tituloconcepto}
  \author{\autor}
  \date{\today}
  \maketitle
\fi


\section{VQE: eigensolver cuántico variacional}

\subsection{Objetivo}
Calcular la energía del estado fundamental de un Hamiltoniano con un algoritmo
híbrido: un circuito cuántico parametrizado prepara estados de prueba y un
optimizador clásico ajusta los parámetros.

\subsection{Conceptos previos}
\begin{concepto}{valor esperado, Hamiltoniano y era NISQ}
Ver capítulos~08 (era NISQ), 12 (Hamiltoniano) y~13 (valores esperados de Pauli).
El \emph{valor esperado} de $H$ en el estado $\ket\psi$ es
$\langle H\rangle=\bra\psi H\ket\psi=\Tr(H\rho)$ con $\rho=\ket\psi\bra\psi$.
\end{concepto}

\subsection{Fundamento teórico}
Los computadores cuánticos actuales tienen ruido y pocos qubits (era NISQ, término
de John Preskill de 2018). En ese contexto son muy usados los algoritmos
\textbf{híbridos cuántico-clásicos}: parte del cálculo corre en un computador
clásico y parte en uno cuántico. El VQE es uno de ellos; su tarea es calcular
la energía aproximada del estado fundamental de un Hamiltoniano. Se aplica
sobre todo en simulación de química cuántica (moléculas como el agua o cadenas de
hidrógeno) y también en optimización y aprendizaje automático cuántico.

\paragraph{Principio variacional.} Para cualquier estado de prueba $\ket\psi$,
\begin{equation}
  E_0\ \le\ \frac{\bra\psi H\ket\psi}{\langle\psi|\psi\rangle} ,
\end{equation}
donde $E_0$ es la energía del estado fundamental. Por eso, minimizar el valor
esperado acerca el resultado a $E_0$.

\paragraph{Los pasos del VQE.}
\begin{enumerate}
  \item \textbf{Hamiltoniano:} se escribe el problema como un Hamiltoniano de
    operadores de Pauli.
  \item \textbf{Estado de prueba:} un circuito parametrizado (rotaciones con
    parámetro y puertas de dos qubits como CNOT o CZ) se aplica a $\ket{0\ldots0}$.
  \item \textbf{Valor esperado:} se mide $\langle H\rangle$ en el computador
    cuántico y se envía al clásico.
  \item \textbf{Optimización:} un método clásico (con o sin gradiente) actualiza
    los parámetros.
  \item Se repiten los pasos 3 y 4 hasta converger o llegar al máximo de iteraciones.
  \item La última iteración da el autoestado y el autovalor óptimos.
\end{enumerate}

\paragraph{Regla del desplazamiento de parámetros.} Para obtener el gradiente
se evalúa el circuito con el parámetro movido:
\begin{equation}
  \frac{dE}{d\theta}=\frac{E(\theta+s)-E(\theta-s)}{2}\quad(s=\pi/2),
  \qquad \theta\leftarrow\theta-\eta\,\frac{dE}{d\theta},
\end{equation}
con $\eta$ la tasa de aprendizaje. Esta forma (con $s=\pi/2$) es exacta para
puertas de rotación $e^{-i\theta P/2}$; con un desplazamiento $\delta$ cualquiera
se obtiene solo una aproximación.

\paragraph{Cálculo de $\langle H\rangle$.} Se puede hacer con el vector de estado,
$\bra\psi H\ket\psi=\sum_{i,j}\bar\psi_iH_{ij}\psi_j$, o con la traza,
$\Tr(H\rho)$. Dan el mismo resultado.

\paragraph{El caso del ejemplo.} El modelo de Heisenberg de dos qubits,
$H=XX+YY+ZZ$, tiene energías $-3$ (estado fundamental, el singlete
$(\ket{01}-\ket{10})/\sqrt2$) y $+1$ (tres veces). El circuito de prueba usa
dos parámetros: $R_y(\theta_1)$ sobre $q_0$, una CNOT $q_0\to q_1$ y
$R_x(\theta_2)$ sobre $q_0$.

\subsection{Código (completamente comentado)}
\lstinputlisting[style=py]{../../fuentes/19_vqe.py}

\subsection{Explicación celda por celda}
\begin{description}
  \item[Celda 2] Define $H=XX+YY+ZZ$ y calcula sus energías exactas por
    diagonalización, solo para comparar.
  \item[Celda 3] Calcula $\langle H\rangle$ de un estado de prueba de las dos formas
    y muestra que coinciden.
  \item[Celda 4] Define el circuito de prueba y la función de energía
    $\Tr(H\rho)$.
  \item[Celda 5] Implementa la regla del desplazamiento y la compara con una
    derivada numérica.
  \item[Celda 6] El ciclo de VQE: 40 iteraciones con tasa de aprendizaje $0.3$.
  \item[Celda 7] Gráfica de convergencia.
\end{description}

\subsection{Salida obtenida}
\lstinputlisting[style=salida]{../../fuentes/salidas/19_vqe.txt}
(El dibujo del circuito no se reproduce aquí.) Es determinista.

\subsection{Interpretación}
Las dos formas de calcular $\langle H\rangle$ dan $-0.2308$, y el gradiente
por desplazamiento coincide con el numérico. Partiendo de
$\theta=(1.3,2.0)$ con $\langle H\rangle=+0.9484$, la energía baja a $-0.51$ en 5
iteraciones, a $-2.996$ en 10 y a $-3.0000$ en 20: el valor exacto del estado
fundamental. Los parámetros terminan en $(-\pi/2,\ \pi)$.

\subsection{Ejercicios}
\begin{enumerate}
  \item Cambia la tasa de aprendizaje a $0.1$ y a $1.0$. ¿Cuántas iteraciones
    hacen falta en cada caso?
  \item Prueba otros parámetros iniciales: ¿converge siempre a $-3$?
  \item Cambia el Hamiltoniano a $H=ZZ$ y observa qué energía mínima se
    obtiene.
\end{enumerate}

\subsection{Errores comunes}
\begin{itemize}
  \item Creer que VQE garantiza el mínimo global: el resultado depende del
    circuito de prueba y de los parámetros iniciales, y puede quedarse en un
    mínimo local.
  \item Olvidar que el valor esperado se estima con muchas mediciones en un
    equipo real; aquí se calcula exacto con el simulador.
\end{itemize}

\ifdefined\enmain\else
  \end{document}
\fi

